\subsection{A precise conditional reduction for the Lefschetz assertions}
\paragraph{J6L (Conditional theorem, proved).}
Let $k$ be algebraically closed, of any characteristic. Let
$j:Y\hookrightarrow X$ be a smooth hyperplane section of a smooth
projective $n$-fold, and use the incidence abelian varieties of J1i.
Fix $\ell\ne\operatorname{char}k$ and put $r=2i-1$. For a smooth
projective $d$-fold $W$, write $L_W$ for cup product by the chosen
hyperplane class. The convention for the Lefschetz lowering
operator is
\[
 \Lambda_W(L_W^a v)=
 \begin{cases}L_W^{a-1}v,&a>0,\\0,&a=0,\end{cases}
 \qquad v\in H^b_{\rm prim}(W,\mathbb Q_\ell).
\]
Thus $\Lambda_W$ lowers degree by two and Tate twist by one; it
has no extra $\mathfrak{sl}_2$ coefficient. Let $B(W)$ mean that
this operator is induced by a rational algebraic correspondence
in $\operatorname{CH}^{d-1}(W\times W)_{\mathbb Q}$.
With this explicit convention the following implications hold:
\begin{enumerate}
\item If $2i\leq n-1$ and $B(X)$ holds, then
$j^*:P^i_{\rm inc}(X)\to P^i_{\rm inc}(Y)$ is an isogeny.
\item If $2i\geq n+1$ and $B(Y)$ holds, then the same homomorphism
is surjective.
\end{enumerate}
Only the relevant degree of $\Lambda_X$, respectively
$\Lambda_Y$, is needed. The cohomological inputs are Deligne's hard Lefschetz theorem
\cite[Theorem 4.1.1, p.~217]{DGJ_Deligne1980} and weak
Lefschetz \cite[(4.1.6), pp.~218--219]{DGJ_Deligne1980}.
The coefficient-free lowering convention is also spelled out in
\cite[Section 4.1, pp.~6--7]{DGJ_Peters2026}.
No integral inverse, equality of
incidence and universal regular equivalence, or middle-degree
fixed-part assertion is part of this conclusion.

\emph{Cohomological detection lemma.}
If $U,V$ are smooth projective, $d=\dim U$, and
$\Gamma\in\operatorname{CH}^{d+b-a}(U\times V)_{\mathbb Q}$
acts as zero on
\[
 H^{2a-1}(U,\mathbb Q_\ell(a))
       \longrightarrow H^{2b-1}(V,\mathbb Q_\ell(b)),
\]
then the induced rational homomorphism
\[
 P^a_{\rm inc}(U)\longrightarrow P^b_{\rm inc}(V)
\]
is zero. Here a rational homomorphism is an element of
$\operatorname{Hom}_{k\text{-av}}(-,-)\otimes\mathbb Q$.
To prove the lemma, set $A=P^a_{\rm inc}(U)$, $e=\dim A$.
Choose a miniversal cycle
$u\in\operatorname{CH}^a(A\times U)$ inducing $[m]$ on $A$,
where $m>0$; see \cite[Lemma 2.3, pp.~104--105]{DGJ_Saito1979}.
Choose a Picard realization of the universal map for $V$, namely
a correspondence $w$ and a closed immersion
$P^b_{\rm inc}(V)\hookrightarrow\operatorname{Pic}^{\rm var}(T)$,
as in \cite[(3.5), p.~110]{DGJ_Saito1979}. The composite
\[
 D=w\circ\Gamma\circ u\in
               \operatorname{CH}^1(A\times T)_{\mathbb Q}
\]
induces the composite of $[m]$, the proposed homomorphism, and
this immersion. Its action on
$H^{2e-1}(A,\mathbb Q_\ell(e))\to
H^1(T,\mathbb Q_\ell(1))$ is zero by the hypothesis.
For a divisor correspondence on $A\times T$, this action is the
Poincaré-dual form of its homomorphism on rational Tate modules.
Consequently the displayed composite homomorphism is zero:
a nonzero homomorphism of abelian varieties has a positive
dimensional abelian-variety image and hence has nonzero rational
Tate-module image. The closed immersion and the invertibility
of $[m]$ in the isogeny category give the lemma. Equivalently one
may clear all denominators first; torsion-free Hom groups give
the same conclusion. The case $A=0$ is immediate.

We now prove the two assertions. The cohomology conventions give
\[
 j_*j^*=L_X,\qquad j^*j_*=L_Y,
 \qquad
 j_*:H^q(Y,\mathbb Q_\ell)
         \longrightarrow H^{q+2}(X,\mathbb Q_\ell(1)).
\]
In the first range, $r\leq n-2$, weak Lefschetz makes $j^*$ an
isomorphism on $H^r$. Every Lefschetz string in $H^r(X)$ has
another step above it, so $\Lambda_X L_X=1$ there. The rational
correspondence
\[
 \beta=\Lambda_X\circ j_*\in
           \operatorname{CH}^{n-1}(Y\times X)_{\mathbb Q}
\]
therefore induces the inverse of $j^*$ on $H^r$:
$\beta j^*=\Lambda_XL_X=1$, and $j^*\beta=1$ follows since
$j^*$ is an isomorphism. The detection lemma applied to the two
composites minus the diagonals proves that $\beta$ is a two-sided
inverse on incidence abelian varieties in the isogeny category.
After clearing denominators, each composite is multiplication
by the same positive integer. Thus the actual restriction
homomorphism is an isogeny.

In the second range, $r\geq n>\dim Y$, every Lefschetz string in
$H^r(Y)$ has a step below it. Hence $L_Y\Lambda_Y=1$ there.
The rational correspondence
\[
 \beta=j_*\circ\Lambda_Y\in
           \operatorname{CH}^{n-1}(Y\times X)_{\mathbb Q}
\]
satisfies $j^*\beta=L_Y\Lambda_Y=1$ on $H^r(Y)$. Detection
therefore gives a right inverse in the isogeny category.
Clearing denominators makes the composite multiplication by a
positive integer on $P^i_{\rm inc}(Y)$, so the actual
restriction is surjective. If $i>\dim Y$, the target is zero
and surjectivity is automatic. This proves J6L.

The proof specifies exactly what cohomological weak Lefschetz
alone lacks: algebraic correspondences for the relevant lowering
operators. It also explains why the stable-range conclusion here
is an isogeny and does not settle the alternate transcription
``isomorphisme''. The middle fixed-factor clause and removal of
the stated algebraicity hypotheses remain active gates.

For comparison, in characteristic zero there is a further precise
reduction relating incidence equivalence to the algebraic part of
the Abel--Jacobi map. Peters' Corollary~7.8, p.~16, states that
if all primitive Lefschetz components of the relevant algebraic
odd cohomology remain algebraic, then these two quotients agree
up to isogeny \cite{DGJ_Peters2026}. Standard conjecture $B(X)$
implies that primitive-component condition; the codimension-two
condition is known, as summarized in his Section~7, pp.~15--16.
This result concerns complex Abel--Jacobi and incidence quotients.
It does not identify the universal regular quotient in all
codimensions and does not remove the positive-characteristic
algebraicity gate in J6L.
